\documentclass[a4paper,10pt]{article}
\newcommand\subdir{/home/koen/LaTeX-setup/}
\input{\subdir/LaTeX-files/imports.tex}
\input{\subdir/LaTeX-files/master-internship.tex}
\usepackage{\subdir/LaTeX-files/aastex}
\usepackage{natbib}
\bibliographystyle{\subdir/LaTeX-files/astroadstitle.bst}

%Set the title, Author and date
\title{Notes Week 29}
\author{Koen Essers}
\tutorial{Master Internship}
\date{\today}

%Set the header style
\header{\tutorialstring}

%Begin the document
\begin{document}
\setuplayout
\section{Thermohaline Mixing}

I simulate a bunch of MS stars with different initial masses. I save profiles depending on the 
central hydrogen fraction to get an even coverage of the \(\mu\) profiles throughout the
MS evolution of the stars. This is important because the \(\mu\) profiles change quite a lot
throughout the MS evolution. We obviously know the age of the MS star when it accretes
mass and thus we can make a pretty good assumation about the mixing.

The problem is then that we could possibly observe the MS stars billions of years after the mass
has been accreted. At this point the \(\mu\) profiles could change and it becomes uncertain.

However, since the abundances are quite uncertain anyway, I think it is okay to just take the
\(\mu\) profile at the time of accretion.

To start, let me plot the \(\mu\) profile from the envelope to the center.
\begin{figure}[ht]
    \marginfignote{0}{Here, we quite clearly see that outer layers of the envelope, which have
    very low density, are not fully ionized. This obviously causes the mean molecular weight to
increase.}
\marginfignote{7}{However, I think this is an unimportant effect. Accreted material will just
the envelope layer further outwards, and the previously unionized material will become ionized.}
\marginfignote{15}{The much more interesting part is close to the core.}
    \centering
    \incplot{w29-mu-env-1}
\end{figure}

Let's look at the profile going from the core outwards.

\begin{figure}[ht]
    \marginfignote{0}{For this particular mass, we see a nearly perfectly diagonal line from the
    mean molecular weight in the core to the mean mean molecular weight in the envelope.}
    \marginfignote{7}{The diagonal line does not shift with the changing core composition, only
    the final height changes. }
    \marginfignote{12}{However, just before the ``diagonal'' line starts, the \(\mu\) profile
    already starts to rise a bit. I think this is pretty interesting, as this seems to change
quite a bit depending on the core composition, and could thus be important for the moment of
accretion.}
    \centering
    \incplot{w29-mu-core-1}
\end{figure}

Below I plot all \(\mu\) profiles for all models, zooming in on this lower portion.

\begin{figure}[ht]
    \centering
    \incplot{w29-mu-core-2}
\end{figure}

Clearly, not only the mass position of the rise in \(\mu\) changes with mass, but also
the slow rise before the diagonal line.

However, I have no intuition for the mean molecular weight of some strongly enriched material,
so before I do a bunch of calculations, let me make a really simple estimate for \(\mu\).
\newsubsection{Simple Derivation}
\begin{questionbox}{}{}
\textbf{\textcolor{keywordscolor}{NOTE:}}
    Simple derivation for myself to refresh how the mean molecular weight is calculated.
\end{questionbox}

The mean molecular weight is defined as:
\begin{equation*}
\mu = \frac{\rho}{nm_{u}}.
\end{equation*}
Here,
\begin{citemize}
    \item  \(\mu\) is the mean molecular weight of a gas.
    \item  \(\rho \) is the mass density of a gas.
    \item  \(n\) is the number density of a gas.
    \item  \(m_{u}\) is the atomic mass unit.
\end{citemize}

For one element, the mass density contributed is
\begin{equation*}
    \rho_{i} = X_{i} \rho.
\end{equation*}
The number density of nuclei is then simply
\begin{equation*}
    n_{i}^\textrm{nuc}  = \frac{\rho_{i}}{A_{i}m_{u}}  .
\end{equation*}

However, if we assume full ionization, we will have 1 nucleus + \(Z_{i} \) electrons.

The ``real'' number density is simply
\begin{equation*}
    n_{i} = (1 + Z_{i})n_{i}^\textrm{nuc}.
\end{equation*}
Now, let me expand and sum over all elements:
\begin{align*}
    n &= \sum_{i} (1 + Z_{i})\frac{\rho X_{i}}{A_{i}m_{u}}\\
      &= \frac{\rho}{m_{u} }\sum_{i} \frac{X_{i}(1 + Z_{i})}{A_{i}}\\
    \frac{nm_{u} }{\rho } &= \sum_{i} \frac{X_{i}(1 + Z_{i})}{A_{i}}\\
    \frac{1}{\mu} &= \sum_{i} \frac{X_{i}(1 + Z_{i})}{A_{i}}\\
\end{align*}

Cool! I am \emph{assuming} full ionization here, which I think is fine. I am pretty sure nearly
the whole MS star is ionized anyway.

Now to compute the mean molecular weight of the accreted gas, we need
\begin{citemize}
    \item  \(X_{i}\), the mass fraction of every element in the accreted gas. We have this.
    \item  \(Z_{i}\), the proton number of every accreted element. We also have this.
    \item  \(A_{i}\), the mass number of every element in the accreted gas. We need to make an
        \emph{assumption} here.
\end{citemize}
Especially for heavy elements, it is quite possible that the element has a few isotopes with
different masses that are all stable and relatively abundant. In this case, the elemental mass
should be the averaged mass of the isotopes.

I think it is fine to take the \href{https://en.wikipedia.org/wiki/Standard_atomic_weight}
{standard atomic weight} in this case. The differences
will be small anyway, especially for the heavy elements.

I will use the package \lstinline[style=output, breaklines=true]|periodictable| written by \citet{kienzle2026} based on data from \citet{meija2013}.

It is then quite straightforward to \(\mu\).

Below I plot the \(\mu\) I calculated from my MESA model with the Monash Abundances.
\begin{figure}[ht]
    \marginfignote{0}{There are a couple of important things here.
\begin{enumd}
    \item  The first elements contribute by far the most to \(\mu\), which is to be expected due
        to their much larger abundances.
    \item  The \(\mu\) computed with the final abundances is larger than for the initial
        envelope abundances.
    \item  \(\mu\) computed from the initial envelope abundances is a bit larger than 
        \(\mu\) in the envelope of the MESA MS models, which have \(\mu \approx 0.60\).
\end{enumd}
    }
    \centering
    \incplot{w28-mu-approx-1}
\end{figure}

I am unsure if the differences in \(\mu\) between the initial models is due to evolution after
the MS, or due to differences between MESA and Monash abundances. But,
in the end, the differences for the predicted abundances in the MS star would probably be small.

The idea is now to use this \(\mu\) and mix the accreted mass, thereby changing \(\mu\) as well,
until the mixed \(\mu\) is equal to the \(\mu\) in the star.
\newsubsection{Simple Derivation 2}
To see how this mixing would work, let me look at the definition of \(\mu\) again.

\begin{equation*}
    \mu = \frac{\rho}{n m_{u} } = \frac{m / V}{Nm_{u} / V} = \frac{m}{Nm_{u}}.
\end{equation*}
Say we mix \(m_{1}\) of gas with mean molecular weight \(\mu_{1} \) with \(m_{2}\) of a gas with \(\mu_{2}\).
We want to know what the resulting \(\mu \) is of the mixed whole.

We would calculate \(N\) using \(\mu \) and \(m\) to obtain:
\begin{equation*}
    N_{i} = \frac{m_{i}}{\mu_{i} m_{u}} \implies \mu_\text{mix} = \frac{m_{1} + m_{2}}{(N_{1}+N_{2})m_{u}} = \frac{m_{1} + m_{2}}{\frac{m_{1}}{\mu_{1}} + \frac{m_{2}}{\mu_{2} } }
\end{equation*}

Or, more conveniently:
\begin{equation*}
    \mu_\text{mix}^{-1}  = \frac{m_{1}\mu_{1}^{-1} + m_{2}\mu_{2}^{-1}  }{m_{1}+m_{2} }\qq{,} m_\text{mix} = m_{1} +m_{2}.
\end{equation*}

I think it should be pretty easy to implement this logic for simulating the dilution of \(\mu \) as the
accreted material starts to mix.

Here is a minimal working example:
\begin{minted}[fontsize=\small]{py}
def effective_mu(profile, M_acc, mu_acc):
    ms = -np.diff(profile.mass)
    mus = (profile.mu[1:] + profile.mu[:-1]) / 2
    arg = np.argmin(mus)
    mus[:arg] = mus[arg]
    mu_current = mu_acc
    m_current = M_acc

    mu_profile = [mu_acc]

    crossing_index = None
    for m, mu in zip(ms, mus):
        mu_last = mu_current
        mu_current = (m + m_current) / ((m / mu) + (m_current / mu_current))
        m_current = m + m_current
        mu_profile.append(mu_current)

    mus = np.concatenate([[mus[0]], mus])
    mu_profile_original = mus
    return mu_profile, mu_profile_original
\end{minted}

Below I show an example of \(0.5\; M_\odot \) being accreted by a \(1.9\; M_\odot \) MS star at various times throughout the MS.

\begin{figure}[ht]
    \marginfignote{0}{Here, accreted mass gets added to the top of the envelope
    of the original star.}
    \marginfignote{5}{As mixing extends into the envelope, the mean molecular weight
    of the mixed material decreases from the initial accreted mean molecular weight
to lower values until at some point, the mean molecular weight of the mixed material
is \emph{lower} than the mean molecular weight of the envelope of the original star.}
    \marginfignote{18}{This critical point roughly determines how deep mixing
    extends.}
    \centering
    \incplot{w28-test-mu-mixing}
\end{figure}

To me, this seems like it works as expected. Another test would be to see how
mixing works with much less accretion. In this case, mixing should cause the
mean molecular weight to drop much quicker.

\begin{figure}[ht]
    \marginfignote{0}{This plot show the same \(1.9\; M_\odot \) model accreting
    ``only'' 0.025 \(M_\odot \) of strongly enriched mass.}
    \marginfignote{4}{And indeed, we see a much quicker drop-off and
    \emph{less} mixing.}
    \marginfignote{7}{We also see that the age of the MS star does change the mixing
    depth a bit. In both case I have shown here, the mixing depth decreases with
age.}
    \marginfignote{13}{This is a nice property, because by using the \(\mu \)
    profile at the time of accretion, we are making an \emph{optimistic} estimate
for the amount of mixing, and thus a \emph{conservative} estimate for the
abundances.}
    \centering
    \incplot{w28-test-mu-mixing-2}
\end{figure}

Now we are quite easily able to obtain the mixing mass. In combination with the
amount of accreted mass and the mass fraction of the elements in the accreted mass,
we are almost ready to compute the envelope abundances in the main sequence stars.

The last thing we need is the envelope abundances of every element in the envelope
of the MS star. I think this is a bit difficult.

\subsection{Interpolated mean molecular weight calculations.}

The basic idea is very similar to the interpolation of the abundances.
First, I interpolate the \(\mu \) profiles for different ages\sidenote{Behind the
scenes, this is then converted to the central hydrogen fraction to remove the
vastly different main-sequence lifetimes of stars with different masses.}.
Then I interpolate these profiles between the masses.

\subsubsection{Main sequence age to central Hydrogen composition.}

This is just a simple linear interpolation as well.

\subsubsection{Interpolated Age and Mass}

I have now created some code that can quickly generate interpolated profiles.
Below, I show two examples. In the first example, I change the age of a MS star
to see how the \(\mu \) profile changes.

\begin{figure}[ht]
    \marginfignote{0}{We can clearly see that, as the star becomes older, the mean
    molecular weight of the star in the core rises. This is to be expected as 
hydrogen turns into helium.}
\marginfignote{8}{We can also see that the slow sharp rise happens at the same
masses for all ages. However, the slow rise before the sharp rise to the core mean
molecular weight value does depend on age. In the next figure, I zoom in on this.}
    \centering
    \incplot{w29-interpolated-age}
\end{figure}

\begin{figure}[ht]
    \marginfignote{0}{The zoom very clearly shows the sharp boundary between the
    slow and fast rise of the mean molecular weight \emph{and} the effect of age.}
    \marginfignote{8}{This is really the regime that we are interested in, because
    the mean molecular weight of the accreted mass will in all cases be pretty close
to my estimated value of 0.64.}
    \centering
    \incplot{w29-interpolated-age-zoom}
\end{figure}

We can also look different masses at a constant age.

\begin{figure}[ht]
    \marginfignote{0}{Now, we can very clearly see that the mass position of the
    rise in molecular weight changes with mass. This makes perfect sense. It is
really just a proxy for the core mass of the star anyway, and lighter stars will have
lighter cores.}
\marginfignote{10}{Since we are comparing stars of different masses at the same
age, we also see that the maximum \(\mu \) in the core increases with mass too. This
also makes perfect sense, since heavier stars have shorter main-sequence lifetimes,
and thus, at the same age will have burned through more of their core hydrogen.}
    \centering
    \incplot{w29-interpolated-mass}
\end{figure}


\bibliography{master.bib}
\end{document}
